\chapter{Введение}

Проблема абсолютной устойчивости является одной из фундаментальных задач теории управления нелинейными системами. Она возникла в середине XX века в работах А.~И.~Лурье, В.~М.~Попова, Р.~Калмана и других исследователей, которые стремились разработать методы анализа устойчивости систем с неопределёнными нелинейными элементами.

\section{Исторический контекст}

Классическая задача абсолютной устойчивости была сформулирована А.~И.~Лурье в 1944 году и получила название \emph{проблемы Лурье}. Основная идея заключается в том, что в реальных системах управления часто встречаются нелинейные элементы (реле, насыщение, гистерезис), точная математическая модель которых неизвестна или слишком сложна для анализа. Однако можно указать класс допустимых нелинейностей, к которому принадлежит реальный элемент, и гарантировать устойчивость системы для \textbf{всех} нелинейностей из этого класса.

В 1957 году В.~М.~Попов предложил частотный критерий абсолютной устойчивости, который стал одним из наиболее эффективных инструментов анализа. В 1960-х годах Р.~Калман и В.~А.~Якубович установили связь между частотными условиями и существованием функций Ляпунова специального вида, что привело к созданию леммы Калмана--Якубовича--Попова (КЯП).

\section{Основные понятия}

\textbf{Абсолютная устойчивость}~--- это свойство замкнутой системы, состоящей из линейной динамической части и статической нелинейности, быть асимптотически устойчивой в целом (глобально) для \emph{всех} допустимых нелинейностей из заданного класса, без необходимости знать точную форму нелинейности.

Ключевая особенность абсолютной устойчивости заключается в том, что она обеспечивает \emph{робастность} системы по отношению к неопределённости в описании нелинейного элемента. Это особенно важно в практических приложениях, где:
\begin{itemize}
    \item точная модель нелинейности недоступна;
    \item параметры нелинейного элемента могут изменяться в процессе эксплуатации;
    \item требуется гарантировать устойчивость для широкого класса нелинейностей.
\end{itemize}

\section{Области применения}

Проблема абсолютной устойчивости находит применение в различных областях:
\begin{itemize}
    \item \textbf{Системы управления с ограничениями}: исполнительные механизмы с насыщением, ограничения по мощности.
    \item \textbf{Адаптивные системы}: системы с неопределёнными параметрами, где нелинейность возникает из-за адаптивных законов.
    \item \textbf{Робастное управление}: анализ устойчивости при неопределённости в описании объекта.
    \item \textbf{Нелинейные наблюдатели}: анализ устойчивости систем с нелинейными корректорами.
    \item \textbf{Скользящие режимы}: связь с теорией скользящих режимов через секторные ограничения.
\end{itemize}

\chapter{Постановка задачи (формулировка Лурье)}

\section{Каноническая форма Лурье}

Рассматривается одномерная (SISO) система в форме Лурье:
\begin{equation}
    \begin{aligned}
        \dot x &= Ax + b\,\varphi(\sigma),\\
        \sigma &= c^\top x, \qquad y = \sigma,
    \end{aligned}
    \label{eq:lurie}
\end{equation}
где:
\begin{itemize}
    \item $x\in\mathbb{R}^n$~--- вектор состояния;
    \item $A\in\mathbb{R}^{n\times n}$~--- матрица линейной части (предполагается, что $A$ гурвицева, т.~е. все собственные значения имеют отрицательные вещественные части);
    \item $b\in\mathbb{R}^n$, $c\in\mathbb{R}^n$~--- векторы входа и выхода;
    \item $\varphi:\mathbb{R}\to\mathbb{R}$~--- статическая нелинейность (не зависит явно от времени);
    \item $\sigma\in\mathbb{R}$~--- вход нелинейности (линейная комбинация состояний);
    \item $y\in\mathbb{R}$~--- выход системы.
\end{itemize}

Передаточная функция линейной части от входа нелинейности $\varphi(\sigma)$ до выхода $\sigma$ имеет вид:
\begin{equation}
    G(s) = c^\top (sI - A)^{-1} b = \frac{N(s)}{D(s)},
\end{equation}
где $N(s)$, $D(s)$~--- полиномы с вещественными коэффициентами, причём степень $N(s)$ не превышает степень $D(s)$.

\section{Определение абсолютной устойчивости}

\begin{definition}[Абсолютная устойчивость]
Система \eqref{eq:lurie} называется \textbf{абсолютно устойчивой} в классе нелинейностей $\Phi$, если:
\begin{enumerate}
    \item нулевое решение $x=0$ является глобально асимптотически устойчивым для \emph{всех} функций $\varphi\in\Phi$;
    \item для любых начальных условий $x(0)\in\mathbb{R}^n$ решение $x(t)$ системы \eqref{eq:lurie} стремится к нулю при $t\to\infty$.
\end{enumerate}
\end{definition}

Важно отметить, что абсолютная устойчивость является более сильным требованием, чем обычная устойчивость для конкретной нелинейности. Она гарантирует устойчивость \emph{априори}, без знания точной формы $\varphi(\cdot)$.

\section{Типы устойчивости}

В контексте абсолютной устойчивости различают:
\begin{itemize}
    \item \textbf{Абсолютная устойчивость в целом} (глобальная): устойчивость для всех начальных условий $x(0)\in\mathbb{R}^n$.
    \item \textbf{Абсолютная устойчивость в малом} (локальная): устойчивость для начальных условий из некоторой окрестности нуля.
    \item \textbf{Абсолютная устойчивость в конечном времени}: достижение нуля за конечное время (связано с конечновременными регуляторами).
\end{itemize}

В классической постановке Лурье рассматривается глобальная абсолютная устойчивость.

\chapter{Классы нелинейностей}

Для формулировки задачи абсолютной устойчивости необходимо задать класс допустимых нелинейностей $\Phi$. Различные классы приводят к различным критериям устойчивости.

\section{Секторные ограничения}

Наиболее распространённым является класс нелинейностей, удовлетворяющих \textbf{секторным ограничениям}.

\begin{definition}[Сектор $[k_1, k_2]$]
Функция $\varphi:\mathbb{R}\to\mathbb{R}$ принадлежит сектору $[k_1, k_2]$, где $0\le k_1 < k_2$, если:
\begin{equation}
    k_1\sigma^2 \le \varphi(\sigma)\sigma \le k_2\sigma^2, \quad \forall \sigma\in\mathbb{R}, \quad \varphi(0)=0.
    \label{eq:sector}
\end{equation}
\end{definition}

Геометрически это означает, что график функции $\varphi(\sigma)$ лежит в секторе между прямыми $y=k_1\sigma$ и $y=k_2\sigma$ (рис.~\ref{fig:sector}).

Частный случай~--- сектор $[0, k]$:
\begin{equation}
    0 \le \varphi(\sigma)\sigma \le k\sigma^2, \quad \forall \sigma\in\mathbb{R}, \quad \varphi(0)=0.
\end{equation}

Этот класс включает:
\begin{itemize}
    \item \textbf{Насыщение}: $\varphi(\sigma) = \mathrm{sat}(\sigma) = \begin{cases} \sigma, & |\sigma|\le 1 \\ \mathrm{sign}(\sigma), & |\sigma|>1 \end{cases}$ (сектор $[0,1]$).
    \item \textbf{Реле}: $\varphi(\sigma) = k\,\mathrm{sign}(\sigma)$ (сектор $[0,k]$).
    \item \textbf{Кубическая нелинейность}: $\varphi(\sigma) = \sigma^3$ (сектор $[0,+\infty)$ при ограниченном $\sigma$).
\end{itemize}

\section{Ограничения на производную}

\begin{definition}[Ограниченная производная]
Функция $\varphi:\mathbb{R}\to\mathbb{R}$ имеет ограниченную производную, если существует $\mu>0$ такое, что:
\begin{equation}
    \left|\frac{d\varphi(\sigma)}{d\sigma}\right| \le \mu, \quad \forall \sigma\in\mathbb{R}.
\end{equation}
\end{definition}

Это условие используется в модифицированных критериях Попова и в анализе систем с супер-твистинг алгоритмом.

\section{Монотонность и глобальная Липшицевость}

\begin{definition}[Монотонная нелинейность]
Функция $\varphi:\mathbb{R}\to\mathbb{R}$ называется монотонной в секторе $[k_1, k_2]$, если:
\begin{equation}
    k_1 \le \frac{\varphi(\sigma_1) - \varphi(\sigma_2)}{\sigma_1 - \sigma_2} \le k_2, \quad \forall \sigma_1\ne\sigma_2.
\end{equation}
\end{definition}

Это условие сильнее секторного ограничения и используется в некоторых расширенных критериях.

\section{Энергетические ограничения}

\begin{definition}[Ограниченная энергия]
Функция $\varphi$ удовлетворяет условию ограниченной энергии, если:
\begin{equation}
    \int_0^\infty \varphi(\sigma(t))\sigma(t)\,dt < \infty
\end{equation}
для всех траекторий системы.
\end{definition}

Это условие используется в частотных критериях, основанных на пассивности.

\section{Связь между классами}

Различные классы нелинейностей связаны между собой:
\begin{itemize}
    \item Монотонная функция в секторе $[k_1, k_2]$ автоматически принадлежит сектору $[k_1, k_2]$.
    \item Функция с ограниченной производной $|\varphi'|\le\mu$ принадлежит сектору $[-\mu, \mu]$.
    \item Секторные ограничения являются наиболее общими и часто используются в практических приложениях.
\end{itemize}

\chapter{Критерии абсолютной устойчивости}

Существует несколько эквивалентных формулировок критериев абсолютной устойчивости: частотные (круговой критерий, критерий Попова) и алгебраические (лемма Калмана--Якубовича--Попова в форме LMI). Все они основаны на построении функций Ляпунова специального вида.

\section{Функции Ляпунова для абсолютной устойчивости}

Классический подход к доказательству абсолютной устойчивости основан на использовании функций Ляпунова вида:
\begin{equation}
    V(x) = x^\top P x + q \int_0^{c^\top x} \varphi(\tau)\,d\tau,
    \label{eq:lyapunov}
\end{equation}
где $P=P^\top>0$~--- положительно определённая матрица, $q\ge 0$~--- неотрицательный параметр.

Производная функции Ляпунова вдоль траекторий системы \eqref{eq:lurie} имеет вид:
\begin{equation}
    \dot V = x^\top(A^\top P + PA)x + 2x^\top P b \varphi(\sigma) + q \varphi(\sigma) \dot\sigma.
\end{equation}

Условие $\dot V < 0$ для всех допустимых нелинейностей приводит к критериям абсолютной устойчивости.

\section{Круговой критерий}

\begin{theorem}[Круговой критерий]
Система \eqref{eq:lurie} абсолютно устойчива в секторе $[0, k]$, если:
\begin{enumerate}
    \item Матрица $A$ гурвицева (линейная часть устойчива).
    \item Передаточная функция $G(s) = c^\top(sI-A)^{-1}b$ не имеет полюсов на мнимой оси.
    \item Найквист-диаграмма $G(j\omega)$ не охватывает и не пересекает окружность с центром в точке $(-1/k, 0)$ и радиусом $1/k$ (или, эквивалентно, точка $(-1/k, 0)$ не лежит внутри замкнутой кривой Найквиста).
\end{enumerate}
\end{theorem}

Геометрическая интерпретация: на комплексной плоскости строится \emph{запрещённая область}~--- круг с центром $(-1/k, 0)$ и радиусом $1/k$. Если кривая Найквиста $G(j\omega)$ целиком лежит вне этого круга, система абсолютно устойчива.

Для сектора $[k_1, k_2]$ запрещённая область представляет собой круг, проходящий через точки $(-1/k_1, 0)$ и $(-1/k_2, 0)$.

\subsection{Связь с положительной действительностью}

Круговой критерий эквивалентен условию строгой положительной действительности передаточной функции:
\begin{equation}
    \Re\{G(j\omega)\} > -\frac{1}{k}, \quad \forall \omega\in\mathbb{R},
\end{equation}
что означает, что Найквист-диаграмма целиком лежит правее вертикальной прямой $x = -1/k$.

\section{Критерий Попова}

Критерий Попова является обобщением кругового критерия и часто даёт менее консервативные результаты.

\begin{theorem}[Критерий Попова]
Система \eqref{eq:lurie} абсолютно устойчива в секторе $[0, k]$, если:
\begin{enumerate}
    \item Матрица $A$ гурвицева.
    \item Существует $\nu\ge 0$ такое, что для всех $\omega\in\mathbb{R}$ выполняется:
    \begin{equation}
        \Re\big\{(1+j\omega\nu) G(j\omega)\big\} > -\frac{1}{k}.
        \label{eq:popov}
    \end{equation}
\end{enumerate}
\end{theorem}

Параметр $\nu$ называется \emph{параметром Попова} и может быть выбран для ослабления условий устойчивости. Геометрически условие \eqref{eq:popov} означает, что модифицированная Найквист-диаграмма $(1+j\omega\nu)G(j\omega)$ целиком лежит правее прямой $x = -1/k$.

\subsection{Преимущества критерия Попова}

Критерий Попова часто менее консервативен, чем круговой критерий, поскольку:
\begin{itemize}
    \item Параметр $\nu$ позволяет учитывать фазовые свойства системы.
    \item Для систем с большим запасом по фазе можно выбрать $\nu>0$ и ослабить условия.
    \item Критерий Попова может гарантировать абсолютную устойчивость даже когда круговой критерий не выполняется.
\end{itemize}

\subsection{Модифицированный критерий Попова}

Для нелинейностей с ограниченной производной $|\varphi'|\le\mu$ можно использовать модифицированный критерий Попова:
\begin{equation}
    \Re\big\{(1+j\omega\nu) G(j\omega)\big\} + \frac{\omega^2\nu^2}{\mu} |G(j\omega)|^2 > -\frac{1}{k}, \quad \forall \omega\in\mathbb{R}.
\end{equation}

\section{Лемма Калмана--Якубовича--Попова (LMI-форма)}

Алгебраическая формулировка критериев абсолютной устойчивости основана на лемме КЯП.

\begin{theorem}[Лемма КЯП]
Пусть $(A, b)$ управляема, $(A, c)$ наблюдаема, и $A$ гурвицева. Система \eqref{eq:lurie} абсолютно устойчива в секторе $[0, k]$ тогда и только тогда, когда существуют матрица $P=P^\top>0$ и скаляр $q\ge 0$ такие, что выполняется линейное матричное неравенство (LMI):
\begin{equation}
    \begin{bmatrix}
        A^\top P + P A & P b - q c\\
        b^\top P - q c^\top & -2q/k
    \end{bmatrix} < 0.
    \label{eq:kyplmi}
\end{equation}
\end{theorem}

\subsection{Преимущества LMI-формулировки}

LMI-формулировка имеет ряд преимуществ:
\begin{itemize}
    \item \textbf{Численная проверка}: LMI могут быть эффективно решены с помощью полуопределённого программирования (SDP).
    \item \textbf{Синтез регуляторов}: LMI позволяют формулировать задачи синтеза как задачи оптимизации.
    \item \textbf{Робастность}: легко учитывать неопределённости в параметрах матриц $A, b, c$.
    \item \textbf{Многомерные системы}: естественным образом обобщается на MIMO-системы.
\end{itemize}

\subsection{Связь с частотными критериями}

LMI \eqref{eq:kyplmi} эквивалентна частотному условию:
\begin{equation}
    \Re\{G(j\omega)\} + q\omega\Im\{G(j\omega)\} > -\frac{1}{k}, \quad \forall \omega\in\mathbb{R},
\end{equation}
где $q = \nu$ в критерии Попова. Таким образом, LMI-формулировка объединяет круговой критерий ($q=0$) и критерий Попова ($q>0$).

\section{Критерий для сектора $[k_1, k_2]$}

Для общего сектора $[k_1, k_2]$ критерии формулируются аналогично:

\textbf{Круговой критерий}: Найквист $G(j\omega)$ не охватывает круг, проходящий через точки $(-1/k_1, 0)$ и $(-1/k_2, 0)$.

\textbf{LMI-критерий}: Существуют $P=P^\top>0$ и $q\ge 0$ такие, что:
\begin{equation}
    \begin{bmatrix}
        A^\top P + P A + k_1 k_2 c c^\top & P b - \frac{k_1+k_2}{2}c - q A^\top c\\
        b^\top P - \frac{k_1+k_2}{2}c^\top - q c^\top A & -q(b^\top c + c^\top b) - \frac{1}{k_2-k_1}
    \end{bmatrix} < 0.
\end{equation}

\chapter{Практический алгоритм проверки абсолютной устойчивости}

\section{Пошаговый алгоритм}

Для проверки абсолютной устойчивости системы рекомендуется следующий алгоритм:

\begin{enumerate}
    \item \textbf{Приведение к форме Лурье}:
    \begin{itemize}
        \item Выделить линейную часть системы: определить матрицы $A, b, c$.
        \item Определить вход нелинейности: $\sigma = c^\top x$.
        \item Записать систему в форме \eqref{eq:lurie}.
    \end{itemize}
    
    \item \textbf{Определение класса нелинейностей}:
    \begin{itemize}
        \item На основе физических свойств нелинейного элемента определить сектор $[k_1, k_2]$ или $[0, k]$.
        \item Если точный сектор неизвестен, использовать консервативную оценку.
    \end{itemize}
    
    \item \textbf{Проверка устойчивости линейной части}:
    \begin{itemize}
        \item Убедиться, что матрица $A$ гурвицева (все собственные значения имеют отрицательные вещественные части).
        \item Если $A$ не гурвицева, абсолютная устойчивость невозможна.
    \end{itemize}
    
    \item \textbf{Выбор критерия проверки}:
    \begin{itemize}
        \item \textbf{Круговой критерий}: удобен для графической проверки, даёт наглядную интерпретацию.
        \item \textbf{Критерий Попова}: часто менее консервативен, требует подбора параметра $\nu$.
        \item \textbf{LMI-критерий}: наиболее удобен для численной проверки и автоматизации.
    \end{itemize}
    
    \item \textbf{Численная проверка}:
    \begin{itemize}
        \item Для частотных критериев: построить Найквист-диаграмму $G(j\omega)$ и проверить геометрические условия.
        \item Для LMI: решить задачу на существование $P>0$ и $q\ge 0$, удовлетворяющих \eqref{eq:kyplmi}.
        \item При необходимости подобрать параметры (например, $\nu$ в критерии Попова или максимальный $k$).
    \end{itemize}
    
    \item \textbf{Верификация моделированием}:
    \begin{itemize}
        \item Провести численное моделирование для репрезентативных нелинейностей из класса (насыщение, релей, кусочно-линейные).
        \item Проверить сходимость траекторий к нулю для различных начальных условий.
    \end{itemize}
\end{enumerate}

\section{Инструменты для проверки}

Для практической реализации алгоритма можно использовать:
\begin{itemize}
    \item \textbf{MATLAB/Simulink}: функции \texttt{nyquist}, \texttt{popov}, пакет \texttt{YALMIP} для решения LMI.
    \item \textbf{Python}: библиотеки \texttt{control}, \texttt{cvxpy} для LMI, \texttt{scipy} для численного моделирования.
    \item \textbf{Специализированные пакеты}: \texttt{SOSTOOLS}, \texttt{ROLMIP} для робастного анализа.
\end{itemize}

\section{Типичные ошибки}

При проверке абсолютной устойчивости следует избегать:
\begin{itemize}
    \item Неправильного определения сектора нелинейности (слишком широкий или узкий).
    \item Проверки только для одной конкретной нелинейности вместо всего класса.
    \item Игнорирования условия гурвицевости матрицы $A$.
    \item Неправильной интерпретации Найквист-диаграммы (учёт полюсов на мнимой оси).
\end{itemize}

\chapter{Детальный пример анализа}

Рассмотрим конкретный пример применения критериев абсолютной устойчивости.

\section{Постановка задачи}

Пусть линейная часть системы описывается передаточной функцией:
\begin{equation}
    G(s) = \frac{1}{s(s+2)} = \frac{1}{s^2 + 2s}.
\end{equation}

В пространстве состояний это соответствует системе:
\begin{equation}
    \begin{aligned}
        \dot x_1 &= x_2,\\
        \dot x_2 &= -2x_2 + \varphi(\sigma),\\
        \sigma &= x_1,
    \end{aligned}
\end{equation}
где матрицы имеют вид:
\begin{equation}
    A = \begin{bmatrix} 0 & 1 \\ 0 & -2 \end{bmatrix}, \quad
    b = \begin{bmatrix} 0 \\ 1 \end{bmatrix}, \quad
    c = \begin{bmatrix} 1 \\ 0 \end{bmatrix}.
\end{equation}

Нелинейность $\varphi(\sigma)$ принадлежит сектору $[0, 1]$, т.~е.:
\begin{equation}
    0 \le \varphi(\sigma)\sigma \le \sigma^2, \quad \forall \sigma\in\mathbb{R}, \quad \varphi(0)=0.
\end{equation}

\section{Проверка условий}

\subsection{Устойчивость линейной части}

Собственные значения матрицы $A$: $\lambda_1 = 0$, $\lambda_2 = -2$. Матрица $A$ не является гурвицевой (имеет нулевое собственное значение), что соответствует интегрирующему звену в передаточной функции. Однако это не препятствует абсолютной устойчивости, если нелинейность обеспечивает стабилизацию.

\subsection{Круговой критерий}

Найквист-диаграмма $G(j\omega)$ для $\omega\in\mathbb{R}$:
\begin{equation}
    G(j\omega) = \frac{1}{j\omega(j\omega+2)} = \frac{1}{-2\omega^2 + j2\omega} = \frac{-2\omega^2 - j2\omega}{4\omega^4 + 4\omega^2}.
\end{equation}

Для кругового критерия необходимо проверить, что кривая Найквиста не охватывает точку $(-1, 0)$ (так как $k=1$). Анализ показывает, что условие выполняется.

\subsection{Критерий Попова}

Для критерия Попова выбираем параметр $\nu = 1{,}2$. Модифицированная передаточная функция:
\begin{equation}
    (1+j\omega\nu)G(j\omega) = (1+j1{,}2\omega)\frac{1}{j\omega(j\omega+2)}.
\end{equation}

Численная проверка показывает, что для всех $\omega\in\mathbb{R}$ выполняется:
\begin{equation}
    \Re\big\{(1+j1{,}2\omega)G(j\omega)\big\} > -1,
\end{equation}
следовательно, система абсолютно устойчива в секторе $[0, 1]$ по критерию Попова.

\subsection{LMI-критерий}

Для проверки LMI-критерия необходимо найти матрицу $P=P^\top>0$ и скаляр $q\ge 0$, удовлетворяющие неравенству \eqref{eq:kyplmi}. Численное решение (например, с помощью \texttt{cvxpy} или \texttt{YALMIP}) подтверждает существование таких $P$ и $q$, что доказывает абсолютную устойчивость.

\section{Интерпретация результатов}

Полученный результат означает, что система будет асимптотически устойчива для \emph{любой} нелинейности из сектора $[0, 1]$, включая:
\begin{itemize}
    \item Насыщение: $\varphi(\sigma) = \mathrm{sat}(\sigma)$.
    \item Релей: $\varphi(\sigma) = \mathrm{sign}(\sigma)$ (с учётом того, что $|\varphi(\sigma)|\le|\sigma|$).
    \item Кусочно-линейные функции.
    \item Любые другие функции, удовлетворяющие секторному ограничению.
\end{itemize}

Это важное свойство, поскольку на практике точная форма нелинейности может быть неизвестна или изменяться в процессе эксплуатации.

\chapter{Численная иллюстрация}
Скрипт \texttt{python/absolute\_stability.py} строит частотные кривые и имитацию замкнутой системы с насыщением $\varphi(\sigma)=\mathrm{sat}(\sigma)$ в секторе $[0,1]$.

\begin{figure}[H]
    \centering
    \includegraphics{nyquist_popov.png}
    \caption{Найквист и кривая Попова для $G(s)=1/(s(s+2))$, сектор $[0,1]$}
\end{figure}

\begin{figure}[H]
    \centering
    \includegraphics{time_response.png}
    \caption{Переходный процесс Лурье-системы с насыщением: стабилизация $x_1,x_2$}
\end{figure}

\chapter{Связь с другими методами теории управления}

\section{Связь с пассивностью}

Абсолютная устойчивость тесно связана с концепцией \emph{пассивности} в теории управления. Система называется пассивной, если для всех входов и начальных условий выполняется:
\begin{equation}
    \int_0^T \sigma(t)\varphi(\sigma(t))\,dt \ge 0, \quad \forall T>0.
\end{equation}

Критерии абсолютной устойчивости можно интерпретировать как условия строгой пассивности модифицированной системы.

\section{Связь с робастным управлением}

Проблема абсолютной устойчивости является частным случаем задачи \emph{робастной устойчивости}, где неопределённость описывается не параметрическими возмущениями, а классом нелинейностей. Методы анализа абсолютной устойчивости находят применение в:
\begin{itemize}
    \item Анализе систем с неопределёнными параметрами.
    \item Синтезе робастных регуляторов.
    \item Оценке запасов устойчивости.
\end{itemize}

\section{Связь со скользящими режимами}

В системах со скользящими режимами нелинейность часто имеет вид $\varphi(\sigma) = k\,\mathrm{sign}(\sigma)$, что соответствует релейной нелинейности в секторе $[0, k]$. Критерии абсолютной устойчивости могут использоваться для анализа устойчивости скользящих режимов и оценки времени достижения скользящей поверхности.

\section{Обобщения и расширения}

Классическая теория абсолютной устойчивости была обобщена на:
\begin{itemize}
    \item \textbf{Многомерные системы} (MIMO): критерии для систем с несколькими входами и выходами.
    \item \textbf{Динамические нелинейности}: нелинейности, зависящие не только от текущего значения, но и от производных.
    \item \textbf{Системы с запаздыванием}: абсолютная устойчивость для систем с временными задержками.
    \item \textbf{Дискретные системы}: критерии для дискретных систем в форме Лурье.
    \item \textbf{Конечновременная устойчивость}: обобщение на системы с конечновременной сходимостью.
\end{itemize}

\chapter{Ограничения и области применимости}

\section{Ограничения классической теории}

Классические критерии абсолютной устойчивости имеют ряд ограничений:
\begin{itemize}
    \item \textbf{Статические нелинейности}: большинство критериев применимы только к статическим нелинейностям $\varphi(\sigma)$, не зависящим явно от времени.
    \item \textbf{Одномерные нелинейности}: классическая теория рассматривает скалярные нелинейности, хотя существуют обобщения на многомерный случай.
    \item \textbf{Консервативность}: критерии могут быть консервативными, т.~е. не гарантировать устойчивость для некоторых систем, которые на самом деле устойчивы.
    \item \textbf{Глобальная устойчивость}: большинство критериев гарантируют глобальную устойчивость, что может быть избыточным требованием.
\end{itemize}

\section{Преодоление ограничений}

Современные исследования направлены на:
\begin{itemize}
    \item Разработку менее консервативных критериев (например, интегральные квадратичные ограничения, IQC).
    \item Учёт динамических нелинейностей и систем с запаздыванием.
    \item Применение методов сумм квадратов (SOS) для нелинейных систем.
    \item Использование машинного обучения для автоматического определения секторов нелинейностей.
\end{itemize}

\chapter{Выводы}

Теория абсолютной устойчивости представляет собой мощный инструмент анализа и синтеза систем управления с неопределёнными нелинейными элементами. Основные выводы:

\begin{enumerate}
    \item \textbf{Робастность}: Абсолютная устойчивость обеспечивает робастность системы по отношению к неопределённости в описании нелинейного элемента, гарантируя устойчивость для всего класса допустимых нелинейностей без необходимости знать их точную форму.
    
    \item \textbf{Многообразие критериев}: Существует несколько эквивалентных формулировок критериев абсолютной устойчивости:
    \begin{itemize}
        \item Частотные критерии (круговой, Попова) дают наглядную геометрическую интерпретацию.
        \item LMI-формулировки (лемма КЯП) обеспечивают эффективную численную проверку и автоматизацию.
    \end{itemize}
    
    \item \textbf{Практическая применимость}: Критерии абсолютной устойчивости находят широкое применение в:
    \begin{itemize}
        \item Анализе систем с ограничениями (насыщение, релей).
        \item Синтезе робастных регуляторов.
        \item Оценке запасов устойчивости.
        \item Анализе адаптивных и нелинейных систем управления.
    \end{itemize}
    
    \item \textbf{Ключевые факторы успеха}:
    \begin{itemize}
        \item Корректное определение сектора нелинейности на основе физических свойств элемента.
        \item Правильный выбор критерия проверки в зависимости от особенностей системы.
        \item Верификация результатов численным моделированием.
    \end{itemize}
    
    \item \textbf{Современные направления}: Теория абсолютной устойчивости продолжает развиваться, включая обобщения на многомерные системы, системы с запаздыванием, динамические нелинейности и применение современных вычислительных методов.
\end{enumerate}

Таким образом, абсолютная устойчивость остаётся актуальной и важной областью теории управления, обеспечивая надёжные методы анализа и синтеза систем с неопределёнными нелинейными элементами.

